\chapter{Infraestructura: despliegue del clúster y servicios base}
\label{cap:infraestructura}

El presente capítulo describe cómo se construye, en la práctica, la base
de infraestructura sobre la que corre el resto del sistema. Se recorren el
clúster Kubernetes local, los servicios centrales que se despliegan sobre
él, el catálogo Iceberg configurado sobre el almacenamiento de objetos y
la ejecución de Spark orquestada por Airflow. Esta última demuestra, de
extremo a extremo, que el mecanismo central de la plataforma funciona
antes de construir ningún flujo de negocio sobre él. El
Capítulo~\ref{cap:estado_arte} justifica cada una de estas decisiones
frente a sus alternativas. Este capítulo describe cómo se llevaron a la
práctica, con los hallazgos operativos reales que surgieron al hacerlo.

\section{Arquitectura de despliegue}
\label{sec:infraestructura_arquitectura}

La plataforma se organiza en una cadena de capas. Un orquestador de
flujos, Airflow, dispara trabajos de cómputo distribuido con Spark sobre
Kubernetes. Estos trabajos leen y escriben contra un almacenamiento de
objetos con formato de tabla, MinIO con Iceberg. Los resultados agregados
se publican en una capa relacional, PostgreSQL. Finalmente, esos
resultados se consumen desde una capa de visualización, Superset, que se
describe en el Capítulo~\ref{cap:visualizacion}. Todos estos servicios se
despliegan como \textit{charts} de Helm~\cite{helm_docs}. La
Figura~\ref{fig:tfm_arquitectura_cadena} resume esta cadena.

\begin{figure}[H]
  \centering
  \resizebox{\textwidth}{!}{%
  \begin{tikzpicture}[
    font=\small,
    node distance=6mm,
    stage/.style={rectangle, draw=azul, thick, rounded corners, fill=gris95,
      text width=2.7cm, align=center, inner sep=5pt, minimum height=1.6cm},
    >=Stealth
  ]
    \node[stage] (airflow) {Airflow\\orquestación};
    \node[stage, right=of airflow] (spark) {Spark\\sobre Kubernetes};
    \node[stage, right=of spark] (minio) {MinIO + Iceberg\\almacenamiento};
    \node[stage, right=of minio] (postgres) {PostgreSQL\\capa \textit{gold}};
    \node[stage, right=of postgres] (superset) {Superset\\visualización};

    \draw[->] (airflow) -- (spark);
    \draw[->] (spark) -- (minio);
    \draw[->] (minio) -- (postgres);
    \draw[->] (postgres) -- (superset);
  \end{tikzpicture}}
  \caption{Cadena de capas de la plataforma.}
  \label{fig:tfm_arquitectura_cadena}
\end{figure}

El clúster subyacente es Kubernetes~\cite{cncf_kubernetes}. Corre en un
único nodo local, con un espacio de nombres dedicado llamado
\texttt{tfm-lakehouse}. La Tabla~\ref{tab:tfm_servicios_desplegados}
resume los servicios desplegados y su papel en esta cadena.

\begin{table}[H]
  \centering
  \small
  \renewcommand{\arraystretch}{1.2}
  \begin{tabular}{p{0.16\textwidth} p{0.26\textwidth} p{0.14\textwidth} p{0.34\textwidth}}
    \hline
    \textbf{Servicio} & \textbf{Tecnología} & \textbf{Versión} & \textbf{Rol} \\
    \hline
    Clúster & k3s & v1.36.4+k3s1 & Sustrato Kubernetes de un único nodo \\
    Almacenamiento & MinIO & \textit{chart} 17.0.21 & Objetos S3, cubo \texttt{lakehouse} \\
    Base relacional & PostgreSQL & \textit{chart} 18.11.3 & Capa \textit{gold} publicada \\
    Orquestación & Airflow & \textit{chart} 1.22.0 & Flujos ETL, ejecutor de Kubernetes \\
    Catálogo de tablas & Iceberg & 1.11.0, \textit{runtime} Spark 3.5 & Catálogo Hadoop sobre MinIO \\
    Procesamiento & Spark & 3.5.9 & Trabajos distribuidos lanzados desde Airflow \\
    \hline
  \end{tabular}
  \caption{Servicios desplegados en el clúster local.}
  \label{tab:tfm_servicios_desplegados}
\end{table}

\section{Clúster Kubernetes local}
\label{sec:infraestructura_k3s}

El clúster se instaló como distribución k3s, en modo de un único nodo. Se
configuró con el modo de red \texttt{host-gw} en lugar del
\textit{backend} VXLAN por defecto, que evita un problema conocido de
encapsulado UDP bajo el entorno de desarrollo WSL2 empleado. k3s se
instala como servicio \texttt{systemd}, lo que exige que dicho gestor de
servicios esté habilitado en el entorno. Bajo WSL2 esto no viene activado
por defecto y requiere un paso de configuración explícito antes de
instalar el clúster.

Sobre el nodo así preparado se creó el espacio de nombres dedicado
\texttt{tfm-lakehouse}. También se añadieron los repositorios oficiales de
Helm de Apache Airflow y Apache Superset. Los \textit{charts} de MinIO y
PostgreSQL, en cambio, se consumen por referencia OCI directa en lugar de
mediante un repositorio Helm clásico, tal y como se detalla en la
Sección~\ref{sec:infraestructura_servicios}.

\section{Servicios base: MinIO, PostgreSQL y Airflow}
\label{sec:infraestructura_servicios}

Los tres servicios centrales se despliegan mediante Helm sobre el clúster.
MinIO se configura con un único cubo llamado \texttt{lakehouse}, con un
único prefijo de aterrizaje de ficheros crudos, \texttt{bronze}. La
Sección~\ref{sec:infraestructura_iceberg} explica por qué el cubo no
necesita ningún otro prefijo. PostgreSQL~\cite{postgresql_docs} se despliega
como una instancia dedicada a la futura capa \textit{gold}. Esta instancia
se mantiene
deliberadamente separada de la base de datos de metadatos que el propio
\textit{chart} de Airflow ya incorpora de forma integrada, para no acoplar
dos sistemas con ciclos de vida distintos. Airflow se despliega con el
ejecutor de Kubernetes, \texttt{KubernetesExecutor}, y con persistencia de
sus DAG en un volumen propio. Los ficheros de flujo se entregan
copiándolos directamente a dicho volumen, en lugar de mediante
sincronización continua con un repositorio Git.

Durante el despliegue se constató un hallazgo operativo real, no
anticipado en el planteamiento inicial. El catálogo gratuito de Bitnami,
proveedor de los \textit{charts} de MinIO y PostgreSQL empleados, redujo
drásticamente su oferta de imágenes públicas a lo largo de 2025. Todas las
imágenes anteriores al recorte se movieron a un registro congelado,
\texttt{bitnamilegacy}, que sigue siendo público y sin autenticación pero
que no recibe ya ninguna actualización de
seguridad~\cite{bitnami_charts_issue_35164}. Este hallazgo obligó a fijar
explícitamente, para cada imagen de MinIO y PostgreSQL, tanto el registro
como la etiqueta exacta. Cada una se verificó como descargable en el
momento del despliegue, en lugar de asumirse de la documentación de cada
\textit{chart}. Se acepta como una limitación conocida, razonable para un
clúster local sin exposición a Internet, pero que exigiría revisarse antes
de cualquier despliegue real.

\section{Catálogo Iceberg sobre MinIO}
\label{sec:infraestructura_iceberg}

El catálogo Hadoop de Iceberg se configura mediante un único catálogo de
Spark, llamado \texttt{lakehouse}, con su raíz de almacén en
\texttt{s3a://lakehouse/warehouse}. Esta ruta se mantiene deliberadamente
distinta del prefijo \texttt{bronze} que aloja los ficheros crudos
aterrizados directamente. Mezclar ambos bajo el mismo prefijo arriesgaría
que la lógica interna de listado y \textit{commit} de Iceberg tropezase
con objetos ajenos a sus propias tablas. Dentro de ese almacén, las capas
\textit{bronze}, \textit{silver} y \textit{gold} se representan como
espacios de nombres de Iceberg, es decir \texttt{lakehouse.bronze},
\texttt{lakehouse.silver} y \texttt{lakehouse.gold}, no como catálogos
independientes. Esto evita triplicar la superficie de configuración sin
necesidad funcional alguna. La Figura~\ref{fig:tfm_bucket_layout} muestra
esta separación dentro del mismo cubo.

\begin{figure}[H]
  \centering
  \resizebox{0.8\textwidth}{!}{%
  \begin{tikzpicture}[
    font=\small,
    level distance=16mm,
    sibling distance=38mm,
    bucket/.style={rectangle, draw=gris50, thick, rounded corners,
      text width=4cm, align=center, inner sep=6pt},
    rawprefix/.style={rectangle, draw=azul, thick, rounded corners,
      fill=gris95, text width=3.6cm, align=center, inner sep=5pt},
    ns/.style={rectangle, draw=azul, thick, rounded corners,
      text width=3.4cm, align=center, inner sep=4pt},
    edge from parent/.style={draw=gris50, thick},
    >=Stealth
  ]
    \node[bucket] {Cubo \texttt{lakehouse}}
      child { node[rawprefix] {\texttt{bronze/}\\aterrizaje crudo} }
      child { node[rawprefix] (wh) {\texttt{warehouse/}\\catálogo Iceberg}
        child { node[ns] {\texttt{lakehouse.bronze}} }
        child { node[ns] {\texttt{lakehouse.silver}} }
        child { node[ns] {\texttt{lakehouse.gold}} }
      };
  \end{tikzpicture}}
  \caption{Separación, dentro del mismo cubo, entre el prefijo de
  aterrizaje crudo y las tablas Iceberg del catálogo.}
  \label{fig:tfm_bucket_layout}
\end{figure}

El cubo se aprovisionó al principio con tres prefijos de aterrizaje,
\texttt{bronze}, \texttt{silver} y \texttt{gold}, a semejanza de esos
mismos tres espacios de nombres. La ingesta del
Capítulo~\ref{cap:ingesta} aterriza en \texttt{bronze} los ficheros crudos
de ORCID y los CVN sintéticos, porque son datos que entran a la
plataforma desde fuera y necesitan un primer punto de aterrizaje antes de
cualquier procesamiento. Las capas \textit{silver} y \textit{gold}, en
cambio, no llegan a aterrizar ningún fichero crudo propio. Se calculan
enteramente dentro de la plataforma, mediante trabajos de Spark que leen
\textit{bronze} y escriben su resultado directamente como tablas Iceberg
bajo el espacio de nombres correspondiente de \texttt{warehouse/}. Al
quedar claro esto durante el diseño de la ingesta y la transformación, los
prefijos \texttt{silver} y \texttt{gold} del cubo se retiraron del diseño
por no cumplir ninguna función, y el cubo quedó con el único prefijo de
aterrizaje que sí se usa, \texttt{bronze}.

La configuración de acceso exigió, además, fijar una cadena de versiones
mutuamente compatibles entre Spark, el \textit{runtime} de Iceberg para
Spark y las bibliotecas de acceso a S3 de Hadoop. Esta configuración
incluye el punto de conexión S3, el direccionamiento por ruta en lugar de
por subdominio, y las credenciales inyectadas por variable de entorno
desde el mismo secreto de Kubernetes que usa MinIO. Una combinación
incorrecta de las versiones de Spark, Iceberg y Hadoop es una fuente bien
documentada de errores de \textit{classpath} difíciles de diagnosticar.
Esta cadena de versiones se verificó de dos formas antes de darla por
buena. Primero, que cada artefacto se pudiera descargar realmente desde su
origen oficial. Después, extrayendo la propia distribución de Spark para
comprobar qué versión de los clientes de Hadoop incluía en la práctica, en
lugar de asumirla de la documentación.

\section{Ejecución de Spark desde Airflow}
\label{sec:infraestructura_spark}

Cada trabajo de Spark se lanza desde una tarea de Airflow que usa
\texttt{KubernetesPodOperator} y ejecuta \texttt{spark-submit} en modo
Kubernetes \textit{client}, contra una imagen de Spark propia que ya
incorpora las dependencias de Iceberg y S3A verificadas en la
Sección~\ref{sec:infraestructura_iceberg}. El propio pod de la tarea de
Airflow actúa así como el proceso controlador de Spark, llamado
\textit{driver}. Este proceso crea y destruye sus propios pods ejecutores
en el mismo espacio de nombres. Para ello se creó una cuenta de servicio
dedicada de Kubernetes con permisos explícitos, limitados a ese espacio de
nombres, en lugar de un permiso de alcance global.

La verificación de este mecanismo, antes de construir ningún flujo de
negocio sobre él, se hizo con un trabajo mínimo. Este trabajo crea un
espacio de nombres y una tabla Iceberg de prueba, inserta unas pocas filas
y las relee. Este primer intento reveló varios problemas reales que no
eran evidentes desde el diseño:

\begin{itemize}
  \item \textbf{Credenciales del controlador.} La configuración de Spark
  para inyectar credenciales en el controlador y en los ejecutores solo
  tiene efecto real sobre los ejecutores en modo \textit{client}. En ese
  modo, el controlador es el propio pod que ya lanzó
  \texttt{spark-submit}, así que esas propiedades se ignoran de forma
  silenciosa sobre él. Las credenciales del controlador se resolvieron
  fijándolas directamente sobre la especificación del propio pod de
  Airflow, no mediante configuración de Spark.
  \item \textbf{Punto de entrada de la imagen.} Sustituir por completo el
  punto de entrada de la imagen oficial de Spark para ejecutar
  \texttt{spark-submit} directamente evita un paso interno de dicha
  imagen que adapta el sistema de usuarios del contenedor al
  identificador de usuario arbitrario con el que se ejecuta el pod. Omitir
  ese paso produce errores de bajo nivel en la resolución del directorio
  de usuario y en el inicio de sesión de Hadoop, que no delatan su causa
  real. La solución consiste en invocar \texttt{spark-submit} como
  argumento del punto de entrada original, no en sustituirlo.
  \item \textbf{Permisos de limpieza.} Los permisos concedidos a la cuenta
  de servicio cubrían la creación y el borrado individual de pods, pero no
  el borrado por lotes que Spark efectúa al finalizar un trabajo para
  limpiar sus propios ejecutores. Sin ese permiso adicional, el trabajo en
  sí terminaba correctamente, pero el pod controlador quedaba en estado de
  error por el fallo de limpieza posterior.
\end{itemize}

\section{Verificación de extremo a extremo}
\label{sec:infraestructura_verificacion}

Una vez corregidos los tres puntos anteriores, la prueba se completó con
dos verificaciones independientes. Este patrón de comprobación se repite
en los capítulos siguientes cada vez que se demuestra un resultado nuevo.
En primer lugar, el propio trabajo debía terminar con éxito y su registro
debía confirmar el número de filas esperado. En segundo lugar, un pod
completamente distinto, sin relación con el trabajo que escribió los
datos, debía poder leer la misma tabla a través del mismo catálogo y
obtener el mismo resultado. Ambas verificaciones se cumplieron. Esto
confirma que el mecanismo central de almacenamiento y procesamiento
distribuido de la plataforma funciona de extremo a extremo, antes de
construir sobre él ningún flujo de ingesta o transformación real.

Esta forma de verificar, y no solo la corrección puntual de cada fallo,
sostiene el objetivo de reproducibilidad de todo el sistema. La
plataforma se diseñó para poder reconstruirse desde cero siguiendo una
guía escrita, y esa guía se valida en la práctica, no solo se redacta,
como se retoma en el Capítulo~\ref{cap:visualizacion}.
